\subsection{反交换代数与交错代数}

定义7. 一个Z型分次A-代数E称为反交换的，如果对所有非零齐次元素x，y\ensuremath{\in}E

\[
x y = (- 1) ^ {\deg (x) \deg (y)} y x.\tag{36}
\]

代数 \(\mathbf{E}\) 称为交错的，如果它是反交换的，并且还满足 \(x^{2} = 0\) （对每个奇次数的齐次元素 \(x\in \mathbf{E}\) ）。

注记. (1) 设 \(\mathbf{E}^{+}\) 为 \(\mathbf{E}\) 的分次子代数，即诸 \(\mathrm{E}_{2n}\) ( \(n \in \mathbf{Z}\) ) 的直和；由定义7可知，若 \(\mathbf{E}\) 是反交换的，则 \(\mathbf{E}^{+}\) 是包含在 \(\mathbf{E}\) 的中心（因而是交换的）中的子代数。

\begin{enumerate}
\def\labelenumi{(\arabic{enumi})}
\setcounter{enumi}{1}
\item
  假设2在E中不是零因子；那么若E是反交换的，则E是交错的，因为对于 \(x \in E\) 是奇次数的齐次元素这一情形，由(36)得 \(x^{2} = -x^{2}\) ，从而 \(2x^{2} = 0\) ，依假设即有 \(x^{2} = 0\) 。
\item
  我们将在 § 7 中详细研究交错代数的重要例子。
\end{enumerate}

引理4. 设E为Z型分次代数，S为一组齐次元素 \(\neq0\) ；E中所有齐次分量 \(x\neq0\) 均对每个 \(y\in S\) 满足关系式(36)的元素组成的集合F是E的分次子代数。

只需注意：(1) 如果 \(x'\) , \(x''\) 是两个次数同为 \(p\) 的齐次元素，\(y\) 是次数为 \(q\) 的齐次元素，并且 \(x'y = (-1)^{pq}yx'\) , \(x''y = (-1)^{pq}yx''\) ，那么也有 \((x' + x'')y = (-1)^{pq}y(x' + x'')\) ； (2) 若 \(x'\) , \(x''\) 是次数分别为 \(p'\) , \(p'', y\) 是次数为 \(q\) 的齐次元素，并且 \(x'y = (-1)^{p'q}yx'\) , \(x''y = (-1)^{p''q}yx''\) ，则

\[
(x ^ {\prime} x ^ {\prime \prime}) y = (- 1) ^ {(p ^ {\prime} + p ^ {\prime \prime}) q} y (x ^ {\prime} x ^ {\prime \prime}).
\]

命题13. 设 E 是类型为 Z 的分次 A-代数，且 S 是由代数 E 的非零齐次元素构成的一个生成系 \(\neq0\) ；为使 E 是反交换的（相应地，交错的），当且仅当 (36) 对所有 \(x\in S\) 和 \(y\in S\) 成立（相应地，该条件成立且进一步有：对所有属于S的奇次齐次元素x有 \(x^{2}=0\) ）。

我们首先考虑反交换代数的情形。根据引理4，由所有齐次分量 \(x \neq 0\) 均对所有 \(y \in S\) 满足 (36) 的元素构成的子代数F包含 S 的所有元素，因此 F = E。若现在 \(F'\) 类似地是 E 的这样的子代数：由所有齐次分量 \(x \neq 0\) 对每个非零齐次元素 \(y \neq 0\) 都满足 (36) 的元素组成，则由上述可知 \(F'\) 包含 S 的所有元素，因此 \(F' = E\) ，这就完成了命题在此情形下的证明。

为证明交错代数情形下的命题，可以假设 E 已经是反交换的；于是立即得出，E 中每个奇次数的齐次元素都具有形式 \(\sum_{i} z_{i} x_{i}\) ，其中 \(z_{i} \in E^{+}\) 且 \(x_{i} \in S\) 是奇次数的（利用 \(E^{+}\) 包含在 E 的中心中这一事实）；由此得出 \(\left(\sum_{i} z_{i} x_{i}\right)^{2} = \sum_{i} z_{i}^{2} x_{i}^{2} + \sum_{i < j} z_{i} z_{j} (x_{i} x_{j} + x_{j} x_{i}) = 0\) ，因为依假设有 \(x_{i}^{2} = 0\) ，且由 (36) 有 \(x_{i} x_{j} + x_{j} x_{i} = 0\) 。

命题 14. 设 E 和 F 是两个 Z 型分次 A-代数，均为反交换（相应地，交错）的。则斜张量积 \(\mathbf{E}^{\mathbf{g}}\otimes_{\mathbf{A}}\mathbf{F}\) （第 7 号）是一个反交换（相应地，交错）代数。

\(E^{g} \otimes_{A} F\) 的一个生成系由诸 \(x \otimes y\) 组成，其中 x（相应地，y）是 E（相应地，F）的齐次元素 \(\neq 0\) 。考虑这样两个元素 \(x \otimes y\) , \(x' \otimes y'\) ，其中 \(\deg(x) = p\) , \(\deg(y) = q\) , \(\deg(x') = p'\) , \(\deg(y') = q'\) ，使得 \(x \otimes y\) 的次数为 \(p + q\) ，而 \(x' \otimes y'\) 的次数为 \(p' + q'\) 。则由定义（第 7 号，公式 (27)）并由 (36)，

\[
(x \otimes y) (x ^ {\prime} \otimes y ^ {\prime}) = (- 1) ^ {q p ^ {\prime}} (x x ^ {\prime}) \otimes (y y ^ {\prime})
\]

\[
\begin{array}{r l} (x ^ {\prime} \otimes y ^ {\prime}) (x \otimes y) & = (- 1) ^ {p q ^ {\prime}} (x ^ {\prime} x) \otimes (y ^ {\prime} y) \\ & = (- 1) ^ {p q ^ {\prime} + p p ^ {\prime} + q q ^ {\prime}} (x x ^ {\prime}) \otimes (y y ^ {\prime}) \end{array}
\]

且命题 13 的判据表明 \(\mathbf{E}^{\mathbf{g}}\otimes_{\mathbf{A}}\mathbf{F}\) 是反交换的，因为 \(pq' + p p' + qq' - q p' \equiv (p + q)(p' + q')\) （模 2）。若进一步 E 和 F 是交错的且 \(p + q\) 是奇数，则数 \(p, q\) 中必有一个是奇数，因此 \((x \otimes y)^2 = \pm (x^2) \otimes (y^2) = 0\) 且命题 13 表明 \(E^g \otimes_A F\) 是交错的。

推论. 设 E 是一个 Z 型反交换（相应地，交错）分次 A-代数。则对每个环同态 \(\rho: \mathrm{A} \to \mathrm{B}\) ，分次 B-代数 \(\rho^{*}(\mathrm{E})\) （第 8 号，注 2）是反交换（相应地，交错）的。

环 B 配以平凡分次可视为一个交错 A-代数，且 \(\rho^{*}(\mathbf{E}) = \mathbf{E}^{\mathfrak{g}}\otimes_{\mathbf{A}}\mathbf{B}\) ，因此可应用命题 14。

注. 设 E 是一个 Z 型反交换分次 A-代数。则由 E 的乘法（§1，第 3 号）定义的从 \(E \otimes_{A} E\) 到 E 的 A-线性映射是分次 A-代数 \(E^{g} \otimes_{A} E\) 到 E 的同态，因为在命题 14 的记号下，在代数 E 中，

\[
(x y) (x ^ {\prime} y ^ {\prime}) = (- 1) ^ {q p ^ {\prime}} (x x ^ {\prime}) (y y ^ {\prime}).
\]

